\documentclass[11pt,letterpaper]{article}
\usepackage[T1]{fontenc}
\usepackage{lmodern,microtype}
\usepackage[margin=.75in,headheight=15pt,headsep=15pt]{geometry}
\usepackage{mathtools,amssymb,booktabs,array,xcolor,enumitem,fancyhdr,hyperref}
\definecolor{ink}{HTML}{17243A}
\definecolor{indigo}{HTML}{293D70}
\definecolor{mint}{HTML}{2A8C82}
\definecolor{muted}{HTML}{647084}
\color{ink}
\hypersetup{colorlinks=true,linkcolor=mint,urlcolor=mint,
 pdfauthor={Igor Rivin},pdftitle={Probability Midterm Study Guide: Through Continuous Distributions}}
\setlength{\parindent}{0pt}
\setlength{\parskip}{5pt}
\setlength{\abovedisplayskip}{7pt}
\setlength{\belowdisplayskip}{7pt}
\clubpenalty=10000
\widowpenalty=10000
\setlist[itemize]{leftmargin=1.4em,itemsep=3pt,topsep=4pt}
\setlist[enumerate]{leftmargin=1.6em,itemsep=6pt,topsep=4pt}
\renewcommand{\P}{\mathbb P}
\newcommand{\E}{\mathbb E}
\DeclareMathOperator{\Var}{Var}
\newcommand{\ind}{\mathbf1}
\newcommand{\topic}[1]{{\Large\bfseries\color{indigo}#1}\par\medskip}
\newcommand{\subtopic}[1]{\par\addvspace{7pt}{\bfseries\color{mint}#1}\par\nobreak}
\newcommand{\checkitem}[1]{\item[$\square$] #1}
\pagestyle{fancy}
\fancyhf{}
\fancyhead[L]{\small\color{muted}Probability \textbar\ Midterm study guide}
\fancyhead[R]{\small\color{muted}Chance, Data, and Decisions}
\fancyfoot[L]{\footnotesize\color{muted}Through continuous distributions}
\fancyfoot[R]{\small\color{mint}\thepage}
\renewcommand{\headrulewidth}{.3pt}
\begin{document}
\topic{Midterm study guide}
{\large Everything through continuous distributions}\par
\textbf{Reading scope.} Chapters 1--5 of \emph{Chance, Data, and Decisions},
including the sections marked ``For later'' on geometric and Poisson laws,
coupon collection, densities, and continuous families. Appendix A supplies
the calculus used here. This guide ends before the joint-distribution chapter;
limit theorems, estimation, testing, and Bayesian inference belong to later units.
Bayes' rule for events is part of this midterm scope.

\subtopic{What you should be able to do without looking up a solution}
\begin{itemize}
\checkitem{Translate a story into outcomes, events, a random variable, and
probabilities; explain any uniformity or independence assumption.}
\checkitem{Count ordered and unordered choices, with and without repetition;
use complements and inclusion--exclusion to avoid overcounting.}
\checkitem{Derive birthday and derangement probabilities rather than guess
from a memorized approximation.}
\checkitem{Use conditional probability, total probability, and Bayes' rule;
distinguish disjointness, pairwise independence, and mutual independence.}
\checkitem{Read or construct a PMF or CDF, normalize a density, and compute
interval probabilities, quantiles, expectations, and variances.}
\checkitem{Use linearity and indicators; know when variances add and when
cross terms are needed. Apply Markov and Chebyshev with their hypotheses.}
\checkitem{Choose a named distribution by its mechanism; state its support,
parameters, mean, and variance, and recognize memorylessness.}
\checkitem{Build coupon-collector stages, find expected remaining waits,
and distinguish a mean collection time from a completion deadline.}
\checkitem{Work with uniform, exponential, and normal laws; standardize with
the standard deviation and use a CDF to transform a variable.}
\end{itemize}

\subtopic{A practical review sequence}
\begin{enumerate}
\item \textbf{Recover the ideas.} Read pages 2--6. Close the guide and write
one example of each mechanism. Explain why a nearby wrong model fails.
\item \textbf{Recover the calculations.} Rework a counting problem,
a conditioning problem, an indicator variance problem, a coupon-stage
problem, and a density problem from the earlier assignments and quizzes.
\item \textbf{Try the mixed practice.} Work page 7 before consulting page 8.
For every error, identify the event or modeling step where it started.
\item \textbf{Use the new assignment for depth.} Problem Set 4 extends the
coupon collector and tail formulas. Its unequal-coupon challenge and
infinite-mean prize are guided enrichment, not extra named laws to memorize.
\end{enumerate}
\textbf{Exam logistics.} This guide describes mathematical scope; it does
not set the exam duration, calculator policy, or allowed formula sheet.

\newpage
\topic{1. Counting and probability spaces}
\subtopic{Choose the experiment before the formula}
For $r$ choices from $n$ objects:
\begin{center}\small
\begin{tabular}{@{}lll@{}}
\toprule
 & Without repetition & With repetition\\
\midrule
Ordered & $n(n-1)\cdots(n-r+1)$ & $n^r$\\[3pt]
Unordered & $\binom nr$ & $\binom{n+r-1}r$\\
\bottomrule
\end{tabular}
\end{center}
The last entry counts multisets from $n$ types, equivalently nonnegative
integer solutions of $x_1+\cdots+x_n=r$. A repeated-letter word with
multiplicities $a_1,\ldots,a_k$ has $(\sum_i a_i)!/\prod_i a_i!$ arrangements.
Divide by a factorial only when changing the order does not change the object.

\textbf{Binomial identities.}
$\binom nr=\binom n{n-r}$,
$\binom nr=\binom{n-1}{r-1}+\binom{n-1}r$,
$(a+b)^n=\sum_{r=0}^n\binom nr a^r b^{n-r}$.
Vandermonde's identity
$\sum_j\binom mj\binom n{r-j}=\binom{m+n}r$
counts a committee by how many members come from each group.

\subtopic{Events, weights, and inclusion--exclusion}
A discrete probability space has outcome weights $w_\omega\ge0$ with
$\sum_\omega w_\omega=1$, and
$\P(A)=\sum_{\omega\in A}w_\omega$.
Only for equally likely outcomes is $\P(A)=|A|/|\Omega|$.
A countably infinite set cannot have a uniform probability giving every
outcome the same weight: zero weights sum to zero, and positive equal
weights have infinite sum.

Know union, intersection, complement, and De Morgan's laws:
$(A\cup B)^c=A^c\cap B^c$ and $(A\cap B)^c=A^c\cup B^c$.
\[
\P(A\cup B)=\P(A)+\P(B)-\P(A\cap B),\qquad
\P\Bigl(\bigcup_i A_i\Bigr)\le\sum_i\P(A_i).
\]
For three events, add single-event probabilities, subtract pair
intersections, then add the triple intersection. For more events, continue
alternating by the size of the intersected set.

\subtopic{Two recurring applications}
\textbf{Birthdays.} With $n$ independent uniform birthdays on $d$ days,
\[
\P(\text{collision})=1-\prod_{k=0}^{n-1}(1-k/d)\quad(n\le d),
\]
and the collision probability is one if $n>d$. For $n\le d$,
\[
1-e^{-n(n-1)/(2d)}\le\P(\text{collision})
\le\min\{1,\tbinom n2/d\}.
\]
The left inequality uses $1-u\le e^{-u}$; the right is a union bound over
pairs. The familiar exponential approximation is close when the accumulated
quadratic terms in $\log(1-u)=-u+O(u^2)$ are small.

\textbf{Derangements.} For a uniform random permutation of $n$ objects,
\[
D_n=n!\sum_{j=0}^n\frac{(-1)^j}{j!},\qquad
\P(\text{exactly }k\text{ fixed points})=\frac{\binom nkD_{n-k}}{n!}.
\]
Set $D_0=1,D_1=0$; $D_n=(n-1)(D_{n-1}+D_{n-2})$ for $n\ge2$.
The no-fixed-point probability approaches $e^{-1}$, with error less than
$1/(n+1)!$. For $n\ge2$, $D_n$ is the nearest integer to $n!/e$.
Exactly $n-1$ fixed points is impossible. The expected number of fixed
points is 1 for $n\ge1$, by indicators.

\newpage
\topic{2. Conditioning, independence, and Bayes}
\subtopic{Three formulas, three jobs}
For $\P(B)>0$,
\[
\P(A\mid B)=\frac{\P(A\cap B)}{\P(B)},\qquad
\P(A\cap B)=\P(B)\P(A\mid B).
\]
Conditioning restricts the possible outcomes to $B$ and rescales their
weights. The order of conditions matters: $\P(A\mid B)$ and
$\P(B\mid A)$ answer different questions.

For a partition $B_1,\ldots,B_r$ with positive probabilities, and $\P(A)>0$
when using Bayes' rule,
\[
\P(A)=\sum_i\P(A\mid B_i)\P(B_i),\qquad
\P(B_j\mid A)=\frac{\P(A\mid B_j)\P(B_j)}
{\sum_i\P(A\mid B_i)\P(B_i)}.
\]
The first is total probability; the second reverses a condition using Bayes.
A tree multiplies probabilities along a path and adds them across disjoint paths.

\subtopic{Worked example: a positive screening test}
Suppose $\P(D)=.01$, $\P(+\mid D)=.9$, and $\P(+\mid D^c)=.05$.
These are stipulated probabilities for a mathematical model.
\[
\P(+)=.9(.01)+.05(.99)=.0585,\qquad
\P(D\mid+)=\frac{.009}{.0585}=\frac2{13}\approx.1538.
\]
A high true-positive rate is not the probability that a positive result
came from a member of group $D$. Both paths to a positive result belong
in the denominator.

\subtopic{Independence must be checked, not inferred from symmetry}
$A,B$ are independent when $\P(A\cap B)=\P(A)\P(B)$.
If $\P(B)>0$, this is equivalent to $\P(A\mid B)=\P(A)$.
Independence is preserved by taking complements.
Two disjoint events of positive probability are dependent: their intersection
has probability zero while the product is positive.

\textbf{Pairwise versus mutual.} Toss two fair coins. Let $A$ mean the first
is heads, $B$ the second is heads, and $C$ that the coins agree.
Each event has probability $1/2$, and each pair intersection has probability
$1/4$. Yet $\P(A\cap B\cap C)=1/4\ne1/8$: the three are pairwise,
but not mutually, independent.

\textbf{A shared hidden coin.} Choose a coin of heads probability $1/4$ or
$3/4$, equally likely, and toss that same coin twice. Then
$\P(H_1)=\P(H_2)=1/2$, but $\P(H_1\cap H_2)=5/16$ and
$\P(H_2\mid H_1)=5/8$. Tosses are independent conditional on the chosen
coin, but dependent when that information is hidden. Equal marginal laws
or symmetry under reordering do not imply independence.

\textbf{Monty Hall.} In the standard model, the host knows the prize location,
always opens a different door containing no prize, and always offers a switch.
Staying wins with probability $1/3$; switching wins with probability $2/3$.
Changing the host's rules changes the conditioning problem. State those rules
before using this answer.

\newpage
\topic{3. Random variables, moments, and bounds}
\subtopic{PMF and CDF}
$X$ assigns a number to each outcome. The event $\{X\ge a\}$ is the set
of original outcomes whose recorded values satisfy that inequality.
$p_X(x)=\P(X=x)$ and $F_X(x)=\P(X\le x)$.
A CDF is nondecreasing and right-continuous, with limits 0 and 1.
\[
\P(a<X\le b)=F_X(b)-F_X(a),\qquad \P(X>a)=1-F_X(a).
\]
At an atom, $\P(X=a)=F_X(a)-F_X(a^-)$, where $F_X(a^-)$ is the left limit.
For integer $X$ and integer $k$,
$\P(X\ge k)=1-F_X(k-1)$.

\subtopic{Weighted averages and spread}
\[
\E[g(X)]=\sum_xg(x)p_X(x),\qquad
\Var(X)=\E[(X-\E[X])^2]=\E[X^2]-\E[X]^2.
\]
For infinite sums, absolute integrability guarantees a finite expectation;
a nonnegative variable can have mean $+\infty$. A valid distribution
need not have finite moments.
\[
\E[aX+bY+c]=a\E[X]+b\E[Y]+c,\qquad
\Var(aX+b)=a^2\Var(X).
\]
Linearity does not require independence. For a sum,
\[
\Var(X+Y)=\Var(X)+\Var(Y)+2\bigl(\E[XY]-\E[X]\E[Y]\bigr).
\]
Independent variables make the last term zero. For $m$ independent
variables with common mean $\mu$ and variance $\sigma^2$, their sum has
mean $m\mu$, variance $m\sigma^2$; their average has variance $\sigma^2/m$.
Variance has squared units; standard deviation has the original units.

\subtopic{Indicators and tail methods}
For $I=\ind_A$, $\E[I]=\P(A)$ and $\Var(I)=\P(A)(1-\P(A))$.
For a count $N=\sum_i\ind_{A_i}$, $\E[N]=\sum_i\P(A_i)$, even with dependence.
For independent uniform birthdays, the expected number of matching
\emph{pairs} is $\binom n2/d$; this is not the collision probability.

For nonnegative integer $N$ and nonnegative real $Y$,
\[
\E[N]=\sum_{m=0}^\infty\P(N>m),\qquad
\E[Y]=\int_0^\infty\P(Y>t)\,dt.
\]
For integrable real $X$ the signed formula is
\[
\boxed{\E[X]=\int_0^\infty(1-F_X(t))\,dt-\int_{-\infty}^0F_X(t)\,dt.}
\]
Adding the two areas instead gives $\E[|X|]$. Apply the nonnegative tail
formula to $X_+=\max(X,0)$ and $X_-=\max(-X,0)$ to see why.
These identities also work for discrete or mixed distributions.

\subtopic{Bounds with hypotheses}
If $Y\ge0$ and $a>0$, $\P(Y\ge a)\le\E[Y]/a$ (Markov).
If $X$ has finite mean $\mu$ and variance $\sigma^2$, then for $t>0$,
$\P(|X-\mu|\ge t)\le\sigma^2/t^2$ (Chebyshev).
Prove Markov using $Y\ge a\ind_{\{Y\ge a\}}$; apply it to
$(X-\mu)^2$ to prove Chebyshev. A bound may be loose or exceed one;
probabilities are always at most one.

\newpage
\topic{4. Discrete laws and coupon collection}
{\small All geometric waits below count trials including the first success.
For $\lambda=0$ the Poisson law is constantly zero.}\par
\begin{center}\footnotesize
\renewcommand{\arraystretch}{1.6}
\begin{tabular}{@{}>{\raggedright\arraybackslash}p{1.18in}p{2.32in}p{.61in}p{1.23in}@{}}
\toprule
Law / support & PMF and mechanism & Mean & Variance\\
\midrule
Uniform on $1,\ldots,n$ & $1/n$; uniform label & $(n+1)/2$ & $(n^2-1)/12$\\
Bernoulli$(p)$; $0,1$ & $p$ at 1, $1-p$ at 0; one indicator & $p$ & $p(1-p)$\\
Binomial$(n,p)$; $0,\ldots,n$ & $\binom nkp^k(1-p)^{n-k}$; independent trials with fixed $n,p$ & $np$ & $np(1-p)$\\
Hypergeometric\newline $(N,K,m)$ & $\binom Kk\binom{N-K}{m-k}/\binom Nm$; $m$ draws without replacement & $mp$ & $mp(1-p)\dfrac{N-m}{N-1}$\\
Geometric$(p)$; $1,2,\ldots$ & $(1-p)^{k-1}p$; first success & $1/p$ & $(1-p)/p^2$\\
Poisson$(\lambda)$; $0,1,\ldots$ & $e^{-\lambda}\lambda^k/k!$; rare-event count & $\lambda$ & $\lambda$\\
\bottomrule
\end{tabular}
\end{center}
For hypergeometric, $p=K/N$, $N>1$, and
$\max(0,m-(N-K))\le k\le\min(m,K)$.
For $N=1$ the count is constant. Bernoulli/binomial allow $0\le p\le1$;
geometric requires $0<p\le1$. If geometric counts \emph{failures}, subtract 1
from its mean and support, leaving variance unchanged.

\textbf{Memorylessness.} For geometric $G$,
$\P(G>s+t\mid G>s)=(1-p)^t=\P(G>t)$ for nonnegative integer $s,t$
when the conditioning event has positive probability.

\textbf{Poisson approximation.} Many independent rare trials with $np\approx\lambda$
can be approximated by $\operatorname{Poisson}(\lambda)$.
$\E[Q(Q-1)]=\lambda^2$ is the quickest route to its variance.
A count of occurrences is not automatically Poisson.

\subtopic{Coupon collection: uniform independent types}
With $k$ distinct types already seen, the next-new-type probability is
$(n-k)/n$ and the stage $W_k$ is geometric with this parameter.
The stage laws remain the same after any past completed history, so the
stages are independent. For $T_n=\sum_{k=0}^{n-1}W_k$,
\[
\E[T_n]=nH_n,\qquad
\Var(T_n)=n^2\sum_{j=1}^n\frac1{j^2}-nH_n.
\]
The expected remaining wait with $j$ types missing is $nH_j$.
The expected time to see $r$ types is $n(H_n-H_{n-r})$.
Integral comparison gives $\log(n+1)\le H_n\le1+\log n$.

For $m\ge1$ draws, missing-type indicators have mean $(1-1/n)^m$, so
\[
\E[M_m]=n(1-1/n)^m,\qquad
\P(T_n>m)\le\min\{1,ne^{-m/n}\}.
\]
They are generally dependent; the stage independence does not imply
missing-type independence. Inclusion--exclusion gives
\[
\P(T_n\le m)=\sum_{j=0}^n(-1)^j\binom nj(1-j/n)^m.
\]
Before any draw, $M_0=n$ and $\P(T_n\le0)=0$. A sufficient $(1-\delta)$
completion deadline is $\lceil n(\log n+\log(1/\delta))\rceil$.
The mean $nH_n$ is not a guarantee. Unequal type probabilities require
a different analysis; the challenge in Problem Set 4 supplies one.

{\small\textbf{Book bonus: Poissonization.} If an independent
$K\sim\operatorname{Poisson}(t)$ determines the draw count, type counts
are independent Poisson$(t/n)$, so the missing count is
Binomial$(n,e^{-t/n})$ and completion has probability $(1-e^{-t/n})^n$.
This product is for a Poisson draw count, not a fixed draw count.}

\newpage
\topic{5. Continuous distributions}
\subtopic{Density, CDF, and quantiles}
A density satisfies $f(x)\ge0$ and $\int_{-\infty}^\infty f(x)\,dx=1$.
\[
F(x)=\int_{-\infty}^x f(t)\,dt,\qquad
\P(a<X<b)=\int_a^bf(t)\,dt,\qquad
\E[g(X)]=\int_{-\infty}^\infty g(x)f(x)\,dx.
\]
The last integral is a finite expectation when $\int |g(x)|f(x)\,dx<\infty$.
For a density-based law, $\P(X=a)=0$, so interval endpoints do not matter.
A density height may exceed one. Where the density is continuous,
$F'=f$; a jump in a CDF is an atom, not an ordinary density value.
A continuous law's median solves $F(m)=1/2$; its $q$-quantile solves $F(x)=q$.

\begin{center}\small
\renewcommand{\arraystretch}{1.55}
\begin{tabular}{@{}p{1.16in}p{3.14in}p{1.18in}@{}}
\toprule
Family & Density and CDF & Mean / variance\\
\midrule
Uniform$(a,b)$ & $f=1/(b-a)$ on $[a,b]$; $F(x)=(x-a)/(b-a)$ inside the interval, 0 below, 1 above & $(a+b)/2$; $(b-a)^2/12$\\
Exponential$(\lambda)$ & $f(x)=\lambda e^{-\lambda x}$, $F(x)=1-e^{-\lambda x}$ for $x\ge0$; both zero for $x<0$ & $1/\lambda$; $1/\lambda^2$\\
Normal$(\mu,\sigma^2)$ & $f(x)=\dfrac{e^{-(x-\mu)^2/(2\sigma^2)}}{\sigma\sqrt{2\pi}}$ on $\mathbb R$; $F(x)=\Phi((x-\mu)/\sigma)$ & $\mu$; $\sigma^2$\\
\bottomrule
\end{tabular}
\end{center}
Here $a<b$, $\lambda>0$, and $\sigma>0$.
\textbf{Units matter.} An exponential \emph{rate} has units reciprocal to time;
a Poisson count parameter for an interval of length $t$ is rate times $t$.
If counts on intervals of length $t$ are Poisson with parameter $\lambda t$,
$\P(\text{next arrival}>t)=\P(\text{zero arrivals})=e^{-\lambda t}$.

\subtopic{Memorylessness and normal standardization}
For exponential $W$ and $s,t\ge0$,
$\P(W>s+t\mid W>s)=e^{-\lambda t}=\P(W>t)$.
Its median is $(\log2)/\lambda$, smaller than its mean.
A uniform wait is not memoryless: its remaining possible interval shrinks.

For $X\sim\operatorname{Normal}(\mu,\sigma^2)$, standardize by
$Z=(X-\mu)/\sigma$, dividing by standard deviation, not variance.
\[
\P(a<X<b)=\Phi\left(\frac{b-\mu}{\sigma}\right)
-\Phi\left(\frac{a-\mu}{\sigma}\right),\qquad \Phi(-z)=1-\Phi(z).
\]
You can leave answers in terms of $\Phi$ when no table is supplied.
An affine transform $aX+b$ is normal with mean $a\mu+b$ and variance
$a^2\sigma^2$ for $a\ne0$; at $a=0$ it is constant.

\subtopic{Transformations and the beta shape mentioned in Chapter 5}
For a monotone transform, rewrite the event before doing calculus.
If $F_X(x)=x^2$ on $[0,1]$ and $Y=X^2$, then
$F_Y(y)=\P(X\le\sqrt y)=y$ on $[0,1]$: $Y$ is uniform.
For $Y=X^2$ on a support crossing zero, use
$\{-\sqrt y\le X\le\sqrt y\}$ for $y\ge0$; both branches matter.

A beta density on $[0,1]$ has shape
$c x^{\alpha-1}(1-x)^{\beta-1}$ with $\alpha,\beta>0$.
Determine $c$ by integration when asked. For example $c x(1-x)$ normalizes
with $c=6$ and has CDF $3x^2-2x^3$ on $[0,1]$.
The general normalizing constant and Bayesian use are later tools.

\newpage
\topic{6. Mixed practice: try before reading the answers}
These questions are a diagnostic, not a prediction of the exam's length or
point weights. Give exact answers or expressions in $\Phi$; show the modeling step.
\begin{enumerate}
\item \textbf{Counting.} A committee of 3 is selected uniformly from 4 women
and 6 men. Find the probability of exactly 2 women and of at least 1 woman.
How many distinguishable arrangements of \texttt{BANANA} are there?
\item \textbf{Birthdays and hats.} Four people have independent uniform
birth months on 12 months. Find the chance of a shared month.
For a random permutation of 4 hats, find the chance of no correct hats,
of exactly 2 correct hats, and the expected number of correct hats.
\item \textbf{Bayes.} Choose a box uniformly. Box A has 2 red and 1 blue
ball; box B has 1 red and 3 blue balls. Draw one ball uniformly from the
chosen box. Given red, what is the probability you chose A?
\item \textbf{Independence.} Roll a fair die. Let $A$ mean even and $B$
mean at least 4. Are they independent? If events $C,D$ are disjoint and
each has positive probability, can they be independent?
\item \textbf{Moments and a CDF.} Let $X=-1,0,2$ with probabilities
$1/4,1/2,1/4$. Find its full CDF, mean, variance, and $\P(X\ge0)$.
For two independent copies, find the mean and variance of their sum.
\item \textbf{Choose a law.} An attempt succeeds independently with
probability $1/4$. Find the chance of exactly 2 successes in 5 attempts;
the chance that the first success is after attempt 3; and the mean and
variance of that waiting time. If $Q\sim\operatorname{Poisson}(2)$, find
$\P(Q\ge2)$ and $\E[Q(Q-1)]$.
\item \textbf{Coupon collection.} An album has 4 equally likely types.
Find its mean completion time and expected remaining wait when 2 types
are missing. Find the completion probability after 5 draws by
inclusion--exclusion. Which calculation needs independence of stages:
adding their means or adding their variances?
\item \textbf{Densities.} A density is $f(x)=c(1+x)$ on $[0,1]$, zero
elsewhere. Find $c$, the full CDF, $\P(X>1/2)$, the mean, variance,
and median. What is $\P(X=1/2)$?
\item \textbf{Continuous laws.} For $W\sim\operatorname{Exp}(2)$, find
$\P(W>3\mid W>1)$ and its median. For $U\sim\operatorname{Uniform}(0,4)$,
find $\P(U>3\mid U>1)$ and its variance.
For $X\sim\operatorname{Normal}(5,4)$, express $\P(3<X<8)$ using $\Phi$.
\item \textbf{Bounds and tails.} A nonnegative variable has mean 3.
Bound $\P(Y\ge12)$. A variable has mean 10 and variance 4;
bound $\P(|X-10|\ge6)$. If $Z$ is uniform on $[-1,2]$, compute
its negative and positive tail areas separately, then find $\E[Z]$
and $\E[|Z|]$.
\end{enumerate}
\subtopic{Final checks before you turn in a solution}
Do probabilities lie in $[0,1]$? Does the PMF sum or the density integrate
to one? Does the CDF include the full support? Is variance nonnegative?
Did you convert time units? Did you explain independence where needed?
For a signed payoff, did you subtract the negative area?

\newpage
\topic{7. Answers and short explanations}
\begin{enumerate}
\item There are $\binom{10}3$ committees. The probabilities are
$\binom42\binom61/\binom{10}3=3/10$ and
$1-\binom63/\binom{10}3=5/6$.
\texttt{BANANA} has $6!/(3!2!)=60$ arrangements.
\item The shared-month probability is
$1-(12\cdot11\cdot10\cdot9)/12^4=41/96$.
$D_4=9$, so no correct hats has probability $9/24=3/8$.
Exactly 2 has probability $\binom42D_2/24=1/4$.
The expected correct count is $4(1/4)=1$; indicators need not be independent.
\item Total red probability is $(1/2)(2/3)+(1/2)(1/4)=11/24$.
Bayes gives $\P(A\mid\text{red})=(1/3)/(11/24)=8/11$.
\item $\P(A)=\P(B)=1/2$, but $\P(A\cap B)=2/6=1/3\ne1/4$.
They are dependent. Disjoint events of positive probability cannot be
independent, because zero intersection differs from their positive product.
\item $F_X(x)=0$ below $-1$, $1/4$ on $[-1,0)$,
$3/4$ on $[0,2)$, and 1 on $[2,\infty)$.
$\E[X]=1/4$, $\E[X^2]=5/4$, $\Var(X)=19/16$,
and $\P(X\ge0)=3/4$. The sum has mean $1/2$ and variance $19/8$.
\item Fixed trial count: $\binom52(1/4)^2(3/4)^3=135/512$.
First-success wait: $\P(G>3)=(3/4)^3=27/64$,
$\E[G]=4$, $\Var(G)=12$.
For Poisson, $\P(Q\ge2)=1-3e^{-2}$ and $\E[Q(Q-1)]=4$.
\item Mean completion: $4H_4=25/3$ draws. Remaining wait:
$4H_2=6$ draws. Completion after 5 draws:
\[
1-4(3/4)^5+6(1/2)^5-4(1/4)^5=15/64.
\]
Adding means uses linearity alone; adding variances needs the independence
of stages (or a separate justification that all cross terms vanish).
\item $c\int_0^1(1+x)\,dx=1$, so $c=2/3$.
$F(x)=0$ below 0, $(2x+x^2)/3$ on $[0,1]$, and 1 above 1.
$\P(X>1/2)=7/12$, $\E[X]=5/9$, $\E[X^2]=7/18$,
$\Var(X)=13/162$.
The median solves $(2m+m^2)/3=1/2$, giving $m=-1+\sqrt{5/2}$.
$\P(X=1/2)=0$.
\item Exponential: $e^{-6}/e^{-2}=e^{-4}$ and median $(\log2)/2$.
Uniform: $(1/4)/(3/4)=1/3$ and variance $4^2/12=4/3$.
Normal: $\sigma=2$, so the answer is $\Phi(3/2)-\Phi(-1)$.
\item Markov gives $3/12=1/4$; Chebyshev gives $4/36=1/9$.
For $Z$, the CDF is $(z+1)/3$ on $[-1,2]$.
The negative area is $\int_{-1}^0(z+1)/3\,dz=1/6$;
the positive area is $\int_0^2(2-z)/3\,dz=2/3$.
Thus $\E[Z]=2/3-1/6=1/2$, while $\E[|Z|]=2/3+1/6=5/6$.
\end{enumerate}
\end{document}
